\section{A normalized Wilson response on a fixed window}
\label{sec:fixed-response}

This section tests a definite response on $0<L<\log2$. It gives
the full arithmetic comparison, a regulated Wilson readout and
its Gaussian first variation, followed by two scoped
obstructions. The calculation does not construct the effective
transfer hypothesized in the exposition.

\subsection{The causal target and its fixed contact term}
\label{subsec:fixed-target}

For $f\in C_c^\infty(I_L)$ let $F$ be its zero extension and
$d=x+L/2$. No nontrivial integer delay acts on this window.
The local constant and pole terms in
\eqref{eq:generator} still give
\begin{align}\label{eq:small-generator}
 (G_{0,L}f)(x)={}&w_0f(x)+
 2\int_0^\infty\nGamma(u)\{f(x)-F(x-u)\}\dd u\notag\\
 &+4\int_0^d\cosh(u/2)f(x-u)\dd u.
\end{align}
This is the nonsymmetric causal generator. The symmetric pole
kernel is $2\cosh((x-y)/2)$, with form
$2|C(f)|^2-2|S(f)|^2$, and must not be substituted for the
causal pole kernel.

The gamma tail outside the available delay interval is local
but does not vanish. Substituting $v=e^{-u/2}$ gives
\begin{equation}\label{eq:gamma-tail}
 T(d):=2\int_d^\infty\nGamma(u)\dd u
 =\log\frac{1+e^{-d/2}}{1-e^{-d/2}}
       +2\arctan(e^{-d/2}).
\end{equation}
Thus the first two terms in \eqref{eq:small-generator} are
$(w_0+T(d))f(x)+2\int_0^d\nGamma(u)\{f(x)-f(x-u)\}\dd u$.
Deleting the tail would change the target.

\begin{proposition}[Full first-response distribution]
\label{prop:contact-response}
Set $A_L=-G_{0,L}$ and, for $u>0$,
\begin{equation}\label{eq:offdiag-response}
 b(u)=2\nGamma(u)-4\cosh(u/2)
     =\frac{2e^{-5u/2}}{1-e^{-2u}}-2e^{u/2}.
\end{equation}
With $c_\epsilon=-w_0-T(\epsilon)$, one has on the smooth core,
in $L^2(I_L)$,
\begin{equation}\label{eq:contact-response}
 (A_Lf)(x)=\lim_{\epsilon\downarrow0}
 \left[c_\epsilon f(x)+
       \int_\epsilon^d b(u)f(x-u)\dd u\right],
\end{equation}
where the integral is zero for $d\le\epsilon$. Moreover
\begin{equation}\label{eq:contact-finite-part}
 c_\epsilon=\log\epsilon+\gamma+\log(2\pi)+O(\epsilon).
\end{equation}
\end{proposition}
\begin{proof}
The gamma part of the bracket is
$-w_0f-2\int_\epsilon^\infty\nGamma(u)(f-F(\,\cdot-u))\dd u$.
Its omitted part tends to zero in $L^2$ because
$\norm{F-F(\,\cdot-u)}_2\le u\norm{F'}_2$ and
$\nGamma(u)=O(1/u)$ near zero. The omitted pole integral
is $O(\epsilon)\norm f_2$ by the convolution bound. Finally,
\eqref{eq:gamma-tail} gives
$T(\epsilon)=-\log\epsilon+\log4+\pi/2+O(\epsilon)$.
Insert the specified
$w_0=-\gamma-\pi/2-3\log2-\log\pi$.
\end{proof}

The finite part in \eqref{eq:contact-finite-part} belongs to
this hard-delay-cutoff convention. A different physical
regulator needs an explicit conversion to it; no freedom to
retune the arithmetic constant is introduced.
The off-diagonal sign also matters. With $z=e^u$,
$b(u)=0$ is $z^3-z-1=0$, whose unique root with $z>1$
gives $u_*=0.2811995743229617\ldots<\log2$.
The response kernel is positive below $u_*$ and negative
above it. This excludes a pointwise nonnegative first-response
kernel on a window longer than $u_*$; it does not exclude
positive Gram operators, whose entries or derivatives need
not be pointwise positive.

\subsection{The lowest-mode cancellation at finite shift}
\label{subsec:absorbed-pole}

The prime-free arithmetic symbol is
\[
 a_0^{<}(p)=\psi\left(\tfrac14+\tfrac p2\right)-\log\pi
                +\frac2{p+1/2}+\frac2{p-1/2}.
\]
The digamma recurrence~\cite[Section~5.5]{DLMFGamma} gives
\begin{equation}\label{eq:absorbed-symbol}
 a_0^{<}(p)=\psi\left(\tfrac54+\tfrac p2\right)-\log\pi
                          +\frac2{p-1/2}.
\end{equation}
Equivalently, with $n_1(u)=e^{-2u}\nGamma(u)$ and $w_1=w_0+4$,
\begin{equation}\label{eq:reduced-generator}
 G_{0,L}=w_1I+
 2\int_0^\infty n_1(u)(I-T_u)\dd u+2R_{-1/2,L},
 \qquad
 (R_{-1/2,L}f)(x)=\int_0^d e^{u/2}f(x-u)\dd u.
\end{equation}
The shift in the local term exactly accounts for removing the
$n=0$ gamma mode. The remaining growing exponential is bounded
as a Volterra kernel on this fixed interval; it is not a
decaying covariance mode.

Absorbing $(p+1/2-\omega)/(p+1/2+\omega)$ into the gamma ratio
in \eqref{eq:factor} yields, for $0<\omega\le1/2$,
\begin{equation}\label{eq:small-transfer-factor}
 K_\omega^{<}(p)=\pi^\omega
 \frac{\Gamma((p+5/2-\omega)/2)}
      {\Gamma((p+5/2+\omega)/2)}
 \left(1-\frac{2\omega}{p-1/2+\omega}\right).
\end{equation}
Its compression equals $V_{\omega,L}$ on $L<\log2$.
It is not the full whole-line zeta quotient. In particular the
apparent pole of this factor must not be interpreted as a pole
of the entire completed-zeta function.

Define
\begin{equation}\label{eq:reduced-beta}
 h_\omega(u)=\frac{2\pi^\omega}{\Gamma(\omega)}
 e^{-(5/2-\omega)u}(1-e^{-2u})^{\omega-1}\mathbf1_{u>0}.
\end{equation}
Inverse Laplace transformation on $\Ree p>1$ gives
\begin{equation}\label{eq:small-transfer-kernel}
 k_\omega^{<}(u)=h_\omega(u)-
 2\omega\int_0^u e^{(1/2-\omega)(u-v)}h_\omega(v)\dd v.
\end{equation}
It is locally integrable at every positive shift, with leading
singularity proportional to $u^{\omega-1}$. Its identity limit
at zero is distributional. Differentiating the corresponding
operators on the smooth core recovers
\eqref{eq:reduced-generator}, including the pole.

\subsection{A regulated readout with identity initial data}
\label{subsec:readout-model}

For a fully specified control, take independent abelian
Gaussian bulk gauge components and six Hermitian scalars with
equal Feynman-gauge covariance
$g^2/(4\pi^2|X-Y|^2)$. Take an independent complex defect
doublet with covariance $\delta_{mn}/|a-b|$.
This is a fixed-gauge Gaussian probe model, not the complete
interacting gauge/defect action.
An explicit regulator uses heat times $s\ge\epsilon^2$:
\begin{align}
 D_\epsilon(X-Y)&=g^2\int_{\epsilon^2}^\infty
 (4\pi s)^{-2}e^{-|X-Y|^2/(4s)}\dd s,\\
 C_\epsilon(a-b)&=4\pi\int_{\epsilon^2}^\infty
 (4\pi s)^{-3/2}e^{-|a-b|^2/(4s)}\dd s.
\end{align}
The probability measure is normalized and independent of the
deformation. Reflection positivity of this cutoff measure
is not assumed.

Use the tangent map $M$ of Section~\ref{sec:projectors},
with columns $(-e_7,e_8,e_9,e_6)$, and the curves
\[
 X_{r,\eta}(t)=(rt,\eta r\varphi(t)),\qquad
 \varphi(t)=t^2(1-t),\quad 0\le t\le1.
\]
The ket runs from $t=1$ to zero, ending at the common reference
$b=0$ with tangent $-e_0$. Let $r=e^x$ and
$\eta(\omega)=\eta_0+\omega$, $\eta_0>0$.
This is a test deformation, not an identification with Loewner
capacity. The color fiber is one-dimensional with its ordinary
metric and an identity contraction at $b$; no extra defect
field is inserted there.

Use the physical pair $q_2,\bar q_2$, and define
\begin{align}
 F_{\omega,\epsilon}(x)&=
 e^{x/2}U_M(C_{e^x,\eta(\omega)})q_{2,\epsilon}(a_{e^x}),\\
 \mathsf K_{\omega,\epsilon}(x,y)&=
 \vev{\Theta_0F_{\omega,\epsilon}(x)F_{\omega,\epsilon}(y)}_\epsilon.
\end{align}
The bra uses reversed ordering and $N=-MR$, as in
Section~\ref{sec:reflection}. Replacing $\bar q_2$ by
the same-charge $\bar q_1$ would give the vanishing endpoint
control of Section~\ref{sec:endpoints}. No common endpoint
supercharge or protected expectation is asserted.

Specify the response on every input by
\begin{equation}\label{eq:retarded-readout}
 (\mathcal V^\epsilon_{\omega,L}f)(x)=f(x)+
 \int_{-L/2}^{x}
 [\mathsf K_{\omega,\epsilon}(x,y)-
  \mathsf K_{0,\epsilon}(x,y)]f(y)\dd y.
\end{equation}
The integral pairs the output row with the source state
$\int_{-L/2}^x F_\omega(y)f(y)\dd y$, minus its baseline.
This supplies a direct identity channel and exact
$\mathcal V^\epsilon_{0,L}=I$, without inverting a compact
Gram operator. Causality in $x$ is imposed by the readout,
not derived from Euclidean reflection. The adjoint in the
fixed unweighted norm is the usual upper-triangular integral
adjoint, with conjugated kernel. It is not contour reversal
with unchanged scalar map.

\subsection{A nonzero but bounded Gaussian first variation}
\label{subsec:gaussian-readout}

The exact fixed-endpoint tangent-connection insertion
\eqref{eq:complex-curvature} applies. Its curvature contains
the scalar and polarization terms of
\eqref{eq:insertion}. The bulge displacement vanishes at
both endpoints; their polarizations and radial weights are
fixed. Both bra and ket vary. More generally one would have
\[
 \partial_\omega\vev{\mathcal O}_\omega
 =\vev{\partial_\omega\mathcal O}_\omega
 -\vev{\mathcal O\,\partial_\omega S_\omega}_{\omega,c},
\]
as well as derivatives of external normalizations or source
identifications. The second term is zero in this model
because its normalized measure is fixed.

At free endpoint order the kernel is
$G_o(x-y)=1/(2\cosh((x-y)/2))$, independent of the bulge,
and the response is $I$. The Wilson dressing gives a nonzero
control. Same-branch gauge/scalar contractions cancel.
The ordinary reflected cross contraction gives, in the
continuum abelian Gaussian model,
\begin{equation}\label{eq:gaussian-kernel}
 \mathsf K_\omega(x,y)=G_o(x-y)
 \exp\{\lambda B_{\eta(\omega)}(r,s)\},
 \qquad r=e^x,\ s=e^y,\quad\lambda=\frac{g^2}{4\pi^2},
\end{equation}
where
\begin{equation}\label{eq:unequal-bulge}
 B_\eta(r,s)=2\eta^2rs\int_0^1\!\int_0^1
 \frac{\varphi'(v)\varphi'(t)}
 {(rv+st)^2+\eta^2[r\varphi(v)-s\varphi(t)]^2}
 \dd v\dd t.
\end{equation}
The exponential is exact only for these independent abelian
Gaussian fields. Its first order in $g^2$ is the direct bulk
sector of Section~\ref{sec:response}. It is not the sum of
the interacting endpoint, defect and counterterm diagrams.

Set $A=(rv+st)^2$ and
$H=[r\varphi(v)-s\varphi(t)]^2$. Direct differentiation gives
\begin{equation}\label{eq:unequal-bulge-derivative}
 \partial_\eta B_\eta(r,s)=4\eta rs
 \int_0^1\!\int_0^1
 \frac{\varphi'(v)\varphi'(t)A}{(A+\eta^2H)^2}
 \dd v\dd t.
\end{equation}
Using $|\varphi'(v)|\le2v$ and
$4rsvt\le(rv+st)^2$ gives
\begin{equation}\label{eq:bulge-uniform-bounds}
 |B_\eta(r,s)|\le2\eta^2,\qquad
 |\partial_\eta B_\eta(r,s)|\le4\eta.
\end{equation}
These estimates include the reference corner $v=t=0$.
For the regulator above, its four-dimensional covariance is
a constant times $(1-e^{-\rho^2/(4\epsilon^2)})/\rho^2$.
It and its radial derivative are bounded by constant
multiples of $\rho^{-2}$ and $\rho^{-3}$, uniformly in
$\epsilon$. On compact positive ranges of $r,s,\eta$,
the displacement of the separation is $O(v+t)$.
The differentiated cross integrands are consequently bounded
by a constant times $vt/(v+t)^2$. Dominated convergence
justifies the continuum kernel and its first derivative in
this model. It does not justify removal of interacting
defect or junction counterterms.

At $\omega=0$,
\begin{equation}
 J_0(x,y)=\lambda G_o(x-y)e^{\lambda B_{\eta_0}(r,s)}
                         \partial_\eta B_{\eta_0}(r,s),
 \qquad
 |J_0|\le2\lambda\eta_0e^{2\lambda\eta_0^2}.
\end{equation}
Uniform first differentiability on the finite square therefore gives
\begin{equation}\label{eq:bounded-readout-variation}
 \mathcal V_{\omega,L}=I+\omega\mathcal J_L+o(\omega)
 \quad\hbox{in operator norm},\qquad
 \norm{\mathcal J_L}\le2L\lambda\eta_0e^{2\lambda\eta_0^2},
\end{equation}
where $\mathcal J_L$ is the lower-triangular integral
operator with kernel $J_0$.
For example,
$B_{0.3}(1,1)=0.0102220294386856\ldots$ and
$\partial_\eta B_{0.3}(1,1)=0.0680705213083742\ldots$.
The derivative is nonzero but finite at $x=y$.

\subsection{Two obstructions and their precise scope}
\label{subsec:readout-obstructions}

\begin{proposition}[Bounded first derivatives cannot match]
\label{prop:bounded-response}
If a response has first derivative on the smooth core equal
to the restriction of a bounded $L^2(I_L)$ operator, it cannot
equal $-G_{0,L}$ there. In particular this excludes
\eqref{eq:bounded-readout-variation}, even after adding a
bounded local multiplier to its derivative.
\end{proposition}
\begin{proof}
Fix $0\ne\chi\in C_c^\infty(I_L)$ and set
$f_N=e^{\ii Nx}\chi$. From \eqref{eq:weil} and the digamma
asymptotic~\cite[Section~5.11]{DLMFGamma},
\begin{equation}\label{eq:modulation-growth}
 \frac{Q_{0,L}[f_N]}{\norm\chi^2}
 =\log N-\log(2\pi)+o(1).
\end{equation}
Indeed the Fourier transform is translated by $N$; split the
integral into its central region and the rapidly decaying
tails of $\widehat\chi$. The two pole moments tend to zero
by integration by parts. If a bounded $A$ equaled $-G_{0,L}$
on the core, then
$\Ree\ip{f_N}{Af_N}/\norm\chi^2$ would tend to $-\infty$,
contrary to boundedness of $A$.
\end{proof}

This uses the localized form to disprove full operator
equality, not to weaken the matching target. It does not
assign a whole-line unweighted Fourier multiplier to the
growing pole. A finite smooth clock change still has a
bounded derivative. A singular clock or a regulator limit
with unbounded first-derivative norms requires a separate
analysis and is not excluded. In particular smoothness at
each fixed cutoff alone is not a continuum obstruction.

\begin{proposition}[A persistent identity channel cannot match]
\label{prop:compact-response}
For every $\omega>0$, $V_{\omega,L}$ is compact. A response
of the form $I+C_\omega$ with compact $C_\omega$ satisfies
\[
 \norm{I+C_\omega-V_{\omega,L}}\ge1.
\]
The Gaussian readout \eqref{eq:retarded-readout} has this form.
\end{proposition}
\begin{proof}
The kernel \eqref{eq:small-transfer-kernel} is in $L^1(0,L)$.
Approximate it there by bounded kernels. Young's convolution
bound gives convergence in operator norm, while every
approximating operator on the finite square is
Hilbert--Schmidt. Hence $V_{\omega,L}$ is compact.
The difference in the proposition is identity plus compact,
whose essential norm is one. The Gaussian kernel difference
is bounded on the finite square, so its retarded operator
is Hilbert--Schmidt.
\end{proof}

A nonzero scalar direct coefficient gives the lower bound
$|m_\omega|$ instead. More generally a nonzero multiplication
operator on this nonatomic interval is not compact.
There is no contradiction with the arithmetic strong initial
limit: compact operators can converge strongly to $I$,
but not in operator norm. A successful response must supply
the singularly concentrating identity limit rather than
retain a separate regular direct channel.

\subsection{Consequences for storage and continuation}

Neither a nonzero Gaussian variation nor a positive reflected
Gram kernel gives the physical balance law
$\norm f^2=\norm{\mathcal V_{\omega,L}f}^2+
\norm{\mathcal B_{\omega,L}f}^2$.
After transfer matching, that law requires
$\norm{\mathcal B_{\omega,L}f}^2
=2\omega Q_{0,L}[f]+o(\omega)$.
In a direction with $Q_{0,L}[f]>0$, the amplitude is of order
$\sqrt\omega$. A state difference that is $O(\omega)$ in its
positive Hilbert norm cannot supply this first-order defect.
This statement is conditional on that norm regularity.

The cumulative defect need not be monotone in shift; no
positivity of every instantaneous generator is imposed here.
The cross-interval bound using the new slab's output defect
and the signed shift changes in Section~\ref{sec:continuation}
remain separate tasks. Nor does this prime-free test account
for the first new generator delay
$-2(\log2)/\sqrt2\,T_{\log2}$ when $\log2<L<\log3$.
The excluded regular readout and the unconstructed arithmetic
Loewner driver must not be conflated with all possible Wilson
realizations.
